\chapter{Einleitung}
Inhalt dieses ersten Kapitels: Motivation, Hintergrund, Aufbau der Arbeit etc.

\section{Mini-\LaTeX-Einführung}

\LaTeX ist eine Sammlung von \TeX-Makros und im Gegensatz zu anderen Textverarbeitungen (wie z.\,B. Word) kein What-you-see-is-what-you-get-System, sondern vielmehr ein What-you-see-is-what-you-mean-System. Man schreibt einen Quellcode, der von einem \LaTeX-System erst verarbeitet und ansprechend visualisiert werden muss. Trotz einer längeren Einarbeitungszeit lassen sich wissenschaftliche Arbeiten -- insbesondere in der Mathematik und in den Naturwissenschaften -- anfertigen, die auch strengsten typographischen Anforderungen genügen.

Grundsätzlich kann mit fast jedem Texteditor direkt in die entsprechenden Dateien mit der Endung \texttt{.tex} aus dem Ordner \texttt{kapitel} geschrieben werden. Will man allerdings eine wissenschaftliche Arbeit verfassen, die optisch ansprechend ist und den gegebenen Normen für solcherlei Texte entspricht, so empfiehlt es, sich einen \LaTeX-Editor (wie z.B. TeXnicCenter oder TeXStudio) zu verwenden und wenigstens die Grundzüge dieses Textsatzsystems zu kennen. Siehe dazu beispielsweise unter \url{ftp://ftp.fernuni-hagen.de/pub/pdf/urz-broschueren/broschueren/a0279510.pdf}, auf der Seite der deutschsprachigen Anwendervereinigung \TeX e.V. (\url{http://www.dante.de/}), der Seite des KOMA-Script Documentation Project (\url{http://www.komascript.de/}) oder in diversen Foren, wie z.\,B. unter \url{http://www.mrunix.de/forums/forumdisplay.php?f=38}, um nur einige Beispiele zu nennen. Man beachte stets, dass es zwar viele Quellen zu diesem Thema gibt, viele jedoch -- aufgrund von Neuentwicklungen oder Unwissenheit der Autoren -- mit Vorsicht zu genießen sind.

Für alle \LaTeX-Neulinge folgen nun zwei einfache Beispiele für \LaTeX-Umgebungen, die eine Idee des Textsatzsystems vermitteln und den Einstieg erleichtern sollen. Eine Umgebung wird fast immer durch folgenden Rahmen bestimmt.
\begin{lstlisting}
\begin{Umgebung}
...
\end{Umgebung}
\end{lstlisting}
Dabei kann es sich um Tabellen, Grafiken, Programmcode oder einen mathematischen Satz handeln. Hier wird dieses Prinzip anhand zweier elementarer Komponenten wissenschaftlicher Arbeiten erläutert, Grafiken und Tabellen.
\begin{beispiel}[Grafiken in \LaTeX]~\label{bsp:grafikeninlatex}
    \begin{figure}[H]
        \centering
        \includegraphics[height=1cm]{./bilder/tud_logo}
        \caption{Bildunterschrift}
    \end{figure}
    Der entsprechende Code dazu sieht wie folgt aus:
    \begin{lstlisting}
\begin{figure}[H]
    \centering
    \includegraphics[height=1cm]{./bilder/tud_logo}
    \caption{Bildunterschrift}
\end{figure}
    \end{lstlisting}
    Es handelt sich hierbei also um die \texttt{figure}-Umgebung, die den Rahmen für externe Grafiken bildet. Der Befehl \lstinline{\centering} zentriert das eingefügte Bild \texttt{tud\_logo.eps} aus dem Ordner \texttt{bilder}. Der Befehl \lstinline{\caption} wird für die Bildunterschrift benötigt, eine entsprechende Nummerierung und der zugehörige Eintrag in das Abbildungsverzeichnis werden automatisch ergänzt. Die Bilddateien sollten dabei im \texttt{eps}-Format vorliegen und im Ordner \texttt{bilder} abgelegt werden. Um beispielsweise \texttt{jpg}-Bilder einzubinden, müssen weitere Änderungen vorgenommen werden, die es selbst zu recherchieren gilt (Tipp: Man recherchiere die Kompilierungsmethode "'PDFLaTeX"`, beziehungsweise "'PDF$\rightarrow$DVI$\rightarrow$PS$\rightarrow$PDF"`).
\end{beispiel}

\begin{beispiel}[Tabellen in \LaTeX]~
    \begin{table}[H]
        \centering
        \begin{tabular}{llll}
            \toprule
                    & Spalte A  & Spalte B  & ... \\
            \midrule
            Zeile 1 & Zelle A1  & Zelle B1 & ... \\
            Zeile 2 & Zelle A2  & Zelle B2 & ... \\
            \bottomrule
        \end{tabular}
        \caption{Tabellenuntertitel}
    \end{table}
    Der entsprechende Code dazu sieht wie folgt aus:
    \begin{lstlisting}
\begin{table}[H]
    \centering
    \begin{tabular}{llll}
        \toprule
                & Spalte A  & Spalte B  & ... \\
        \midrule
        Zeile 1 & Zelle A1  & Zelle B1 & ... \\
        Zeile 2 & Zelle A2  & Zelle B2 & ... \\
        \bottomrule
    \end{tabular}
    \caption{Tabellenuntertitel}
\end{table}
    \end{lstlisting}
    Dieses Beispiel schachtelt die Umgebungen \texttt{table} und \texttt{tabular}. Wie in Beispiel \ref{bsp:grafikeninlatex} gibt es eine Rahmenumgebung, \texttt{table}, die analog zu \texttt{figure} zu verstehen ist. Innerhalb dieser Umgebung wird nun kein externes Dokument eingebunden, sondern innerhalb von \LaTeX~eine Tabelle mit der \texttt{tabular}-Umgebung erzeugt. Näheres dazu kann den zuvor genannten Quellen entnommen werden.
\end{beispiel}

Ausnahmen dieser "`\lstinline|begin|/\lstinline|end|-Regel"' sind z.\,B. die \$-Umgebung, sie ermöglicht die Verwendung von Formeln innerhalb des Fließtextes, und die \texttt{cite}-Umgebung für Literaturverweise.

\begin{beispiel}[Formeln]~
    Es gilt $a^2+b^2=c^2$. Will man dies hervorheben so schreibt man besser
    \[ a^2+b^2=c^2. \] \newpage
    \begin{lstlisting}
        Es gilt $a^2+b^2=c^2$. Will man dies hervorheben so schreibt man besser
        \[ a^2+b^2=c^2. \]
    \end{lstlisting}
\end{beispiel}

\begin{beispiel}[Zitate in Latex]~
Tipps und Tricks zur Formulierung mathematischer Gedanken findet man z.\,B. in \cite{beutelspacher2004trivial}.
    \begin{lstlisting}
Tipps und Tricks zur [...] in \cite{beutelspacher2004trivial}.
    \end{lstlisting}
    "`beutelspacher2004trivial"' liefert dabei den eindeutigen Schlüssel, der jedem verwendeten Buch in der Datei \texttt{literatur.tex} im Ordner \texttt{literatur} zugeordnet werden muss.

    Üblicherweise arbeitet man heutzutage mit BibTex, ein Programm zur vereinfachten Erstellung von Literaturangaben und "=verzeichnissen. Hier wurde zunächst darauf verzichtet um ohne Umwege eine lauffähige Version bereitzustellen. Falls man mit BibTex arbeiten möchte, so kann auf die Datei \texttt{literatur.bib} zurückgegriffen werden, aber man muss noch weitere Anpassungen vornehmen. Beispielsweise muss der Befehl \lstinline|\begin{thebibliography}{sotief}
	\bibitem[1]{Courbariaux}
	\bibitem{pmlr-v37-ioffe15}
    
\end{thebibliography}| durch die drei darüber stehenden Kommandos in der \texttt{main.tex} ersetzt werden und der Editor entsprechend angepasst werden. Die Recherche dieser Anpassungen ist hier dem engagierten Studenten überlassen.
\end{beispiel}

Man beachte, dass auch die hier gegebenen Beispiele in einer Umgebung stehen. Es handelt sich hierbei um die \texttt{ntheorem}-Umgebung, die insbesondere für mathematische Arbeiten von Bedeutung ist. Welche mathematischen Umgebungen in dieser Vorlage bereits definiert sind und in welcher Form dies geschieht lässt sich im letzten Teil der \texttt{header.tex} nachlesen. Näheres dazu erfährt man in der Dokumentation zur \texttt{ntheorem}-Umgebung.


\section{weitere Unterkapitel}

$\vdots$ 